% !TEX root = ../algebra_02.tex

\section{Собственные значения. Линейная независимость собственных векторов, принадлежащих разным собственным значениям}

\begin{Definition}{Собственное значение и собственный вектор}{}
	Пусть $\alpha\in\End V$. Число $\lambda\in K$ называется собственным значением оператора $\alpha$, если существует ненулевой вектор $v\in V$ такой, что
	\[
	\alpha v=\lambda v
	\]
	Такой вектор $v$ называется собственным вектором, отвечающим собственному значению $\lambda$.
\end{Definition}

\begin{Definition}{Собственное подпространство}{}
	Собственным подпространством, отвечающим $\lambda$, называется
	\[
	V_\lambda=\{v\in V\mid \alpha v=\lambda v\}.
	\]
	Эквивалентно,
	\[
	V_\lambda=\Ker(\alpha-\lambda\operatorname{id}_V).
	\]
\end{Definition}

Фундаментальным свойством собственных векторов является их независимость при различии собственных значений.

\begin{Theorem}{Собственные векторы с разными собственными значениями}{}
	Пусть $\lambda_1,\ldots,\lambda_k$ --- попарно различные собственные значения оператора $\alpha$. Если $v_i$ --- собственный вектор, отвечающий $\lambda_i$ (то есть $v_i\in V_{\lambda_i}$ и $v_i\ne 0$), то векторы $v_1,\ldots,v_k$ линейно независимы.
\end{Theorem}

\begin{proof}
	Докажем индукцией по $k$. При $k=1$ утверждение очевидно, так как $v_1\ne 0$.

	Пусть выполнено равенство:
	\[
	a_1v_1+\cdots+a_kv_k=0
	\].
	Применим к этому равенству оператор $\alpha$:
	\[
	a_1\lambda_1v_1+\cdots+a_k\lambda_kv_k=0
	\].
	Умножим исходное равенство на $\lambda_k$ и вычтем его из последнего:
	\[
	a_1(\lambda_1-\lambda_k)v_1+\cdots+ a_{k-1}(\lambda_{k-1}-\lambda_k)v_{k-1}=0
	\].
	По предположению индукции векторы $v_1,\ldots,v_{k-1}$ линейно независимы. Так как $\lambda_i\ne\lambda_k$ при $i<k$, коэффициенты $(\lambda_i-\lambda_k)$ не равны нулю, откуда получаем $a_1=\cdots=a_{k-1}=0$. Тогда из исходного равенства следует $a_kv_k=0$, а так как $v_k \ne 0$, получаем $a_k=0$.
\end{proof}

\begin{Corollary}{Прямая сумма собственных подпространств}{}
	Если $\lambda_1,\ldots,\lambda_k$ --- различные собственные значения $\alpha$, то
	\[
	V_{\lambda_1}+\cdots+V_{\lambda_k}
	=
	V_{\lambda_1}\oplus\cdots\oplus V_{\lambda_k}.
	\]
\end{Corollary}

\begin{Corollary}{Оценка числа собственных значений}{}
	Пусть $\lambda_1,\ldots,\lambda_k$ --- все различные собственные значения оператора $\alpha$.
	Тогда
	\[
	g_{\lambda_1}+\cdots+g_{\lambda_k}\le \dim V.  
	\]
	($g_{\lambda_i}$  --- геометрическая кратность, опр. в следующем вопросе). В частности, у оператора на $n$-мерном пространстве не больше $n$ различных собственных значений.
\end{Corollary}

\begin{Corollary}{Достаточное условие диагонализируемости}{}
	Если $\dim V=n$ и у оператора $\alpha$ есть $n$ различных собственных значений, то $\alpha$ диагонализируем.
\end{Corollary}

\begin{proof}
	Выберем ненулевой собственный вектор $v_i\in V_{\lambda_i}$ для каждого собственного значения.
	По теореме эти $n$ векторов линейно независимы.
	Значит, они образуют базис $V$, и в этом базисе матрица оператора диагональна.
\end{proof}